\documentclass[11pt]{article}

\usepackage[margin=1in]{geometry}
\usepackage[T1]{fontenc}
\usepackage[utf8]{inputenc}
\usepackage{lmodern}
\usepackage{hyperref}
\usepackage{enumitem}
\usepackage{longtable}
\usepackage{array}
\usepackage{amsmath}
\usepackage{amssymb}

\hypersetup{
  colorlinks=true,
  linkcolor=blue,
  urlcolor=blue
}

\setlist[itemize]{leftmargin=1.4em}
\setlist[enumerate]{leftmargin=1.6em}
\emergencystretch=3em
\renewcommand{\arraystretch}{1.15}

\title{\textbf{NEC Math Companion Syllabus}\\
Mathematical Foundations for the NEC / MoE Thesis}
\author{Companion to \texttt{NEC\_thesis\_syllabus.md}}
\date{Created June 2026}

\begin{document}
\maketitle

\begin{center}
\begin{tabular}{ll}
\textbf{Purpose} & Deepen only the math that supports the NEC thesis \\
\textbf{Orientation} & Connect new ML math to prior physics/math background \\
\textbf{Main split} & Proposal-level math first; thesis-depth math after proposal \\
\end{tabular}
\end{center}

\tableofcontents
\newpage

\section{Verdict and Scope}

You should go deeper in math, but not equally everywhere. The NEC thesis needs a clear mathematical spine:

\begin{enumerate}
  \item Mixture-model probability.
  \item Optimization and training dynamics.
  \item Load balancing and expert collapse.
  \item HMM / Markov-switching math.
  \item Financial validation statistics.
\end{enumerate}

Information theory and ensemble bias-variance are useful supporting material. They should appear if they help explain gates, entropy, regularization, and why multiple experts might help. They should not become the main thesis.

\begin{quote}
Before proposal: understand and implement the math. After proposal: deepen derivations for the thesis and defense.
\end{quote}

\section{Connection to the Current Syllabus}

\begin{longtable}{>{\raggedright\arraybackslash}p{0.28\textwidth} >{\raggedright\arraybackslash}p{0.28\textwidth} >{\raggedright\arraybackslash}p{0.34\textwidth}}
\textbf{Current syllabus module} & \textbf{Math topic} & \textbf{Thesis use} \\
\hline
Module 2: Linear algebra & Matrix calculus, projections, SVD & Regression, PCA, gradients, softmax gate \\
Module 3: Probability/statistics & Conditional probability, likelihood, entropy & MoE probability, losses, gate diagnostics \\
Module 4: Linear/logistic models & Softmax, convex losses, regularization & NEC classifier/gate math \\
Module 5: Validation & Sampling error, purging, multiple testing & Defend empirical claims \\
Module 6: Neural networks & Backprop, optimization & Train corrected NEC \\
Module 7: Sequence models & Recurrences, hidden states & LSTM gate encoder \\
Module 8: Mixture models & Latent variables, EM & Mathematical heart of MoE \\
Module 9: HMMs/regimes & Dynamic programming, Markov chains & Hamilton baseline \\
Module 10: MoE theory & Routing, load balancing & Corrected NEC design \\
Module 11: Empirical asset pricing & Factor math, residual returns & Finance interpretation \\
\end{longtable}

\section{Bridge From Previous Math}

You are not starting from zero. Several pieces of your physics/math background transfer directly:

\begin{itemize}
  \item Linear algebra / tensors -> neural network layers, PCA, covariance matrices, embeddings, matrix calculus.
  \item Calculus and variational principles -> optimization, regularized objectives, Lagrangian penalties, gradient descent.
  \item ODEs / dynamical systems -> state evolution intuition, though HMMs are stochastic state models rather than deterministic systems.
  \item Sturm-Liouville / spectral thinking -> useful for covariance and PCA intuition, but not the center of NEC.
  \item Green's functions / propagators -> useful analogy for forward recursions, but HMM forward-backward is discrete probabilistic recursion.
  \item Quantum mechanics / statistical mechanics -> useful analogy for softmax, partition functions, entropy, and variational free energy.
  \item Probability -> direct foundation for latent-variable models, mixture likelihoods, and regime inference.
\end{itemize}

The new math is less physics-flavored and more statistical ML-flavored: conditional likelihoods, latent variables, gates, EM, regularization, and financial validation.

\section{Depth Map}

\small
\begin{longtable}{>{\raggedright\arraybackslash}p{0.22\textwidth} >{\raggedright\arraybackslash}p{0.24\textwidth} >{\raggedright\arraybackslash}p{0.29\textwidth} >{\raggedright\arraybackslash}p{0.11\textwidth}}
\textbf{Topic} & \textbf{Pre-proposal depth} & \textbf{Thesis-depth after proposal} & \textbf{Priority} \\
\hline
Mixture-model probability & Derive MoE equation and implement soft MoE & Likelihood derivation, EM comparison, log-sum-exp, responsibilities & Highest \\
Optimization & Backprop, Adam, regularization, gradient flow through gates & Nonconvex training, collapse dynamics, penalty methods & Highest \\
Load balancing & Know formula and implement it & Derive auxiliary loss, analyze utilization/entropy tradeoff & Highest \\
HMM / dynamic programming & Implement forward algorithm and understand Hamilton filter & Forward-backward, Viterbi, Baum-Welch, Markov-switching regression & High \\
Financial validation statistics & Implement purging, rank-IC, backtest costs & Multiple testing, deflated Sharpe, uncertainty of IC & High \\
Information theory & Entropy, KL, cross-entropy & Gate entropy, mutual information diagnostics if useful & Medium \\
Bias-variance / ensembles & ESL-level bias-variance & Optional expert-ensemble discussion & Low-medium \\
Stochastic calculus & Skip & Skip unless supervisor requires it & Low \\
Measure theory & Skip & Skip & Low \\
\end{longtable}
\normalsize

\section{Pillar 1: Linear Algebra and Matrix Calculus}

\textbf{Connected current modules:} Module 2, Module 4, Module 6.

\textbf{Why it matters.} MLP experts, LSTM gates, softmax gates, ridge/logistic baselines, PCA, and covariance diagnostics are all linear algebra plus nonlinearities.

\textbf{Previous-knowledge bridge.} Inner products, basis changes, eigenvectors, and quadratic forms already show up in physics. Here the same objects become regression, covariance, and gradient machinery.

\textbf{Readings by section number.}
\begin{itemize}
  \item Strang, \emph{Linear Algebra and Its Applications}: Sec 3.2-3.4, 5.1-5.3, 6.1-6.3.
  \item Matrix Cookbook: Sec 2.1, 2.4, 2.5, 8.1, 8.2.
  \item ESL: Ch 3.2-3.4.
\end{itemize}

\textbf{Focused exercise set.}
\begin{itemize}
  \item Strang Sec 3.3: Exercises 1, 3, and 7.
  \item Strang Sec 6.3: Exercises 1 and 3.
  \item ESL Ch 3: Exercises 3.4 and 3.6.
  \item Applied: derive and implement ridge regression from scratch.
  \item Applied: derive the softmax Jacobian and verify it numerically.
\end{itemize}

\textbf{Thesis deliverable.} You should be able to write
\[
\pi(x) = \mathrm{softmax}(Wh(x)+b)
\]
and explain how gradients flow from the prediction loss through \(\pi(x)\) into the gate parameters.

\section{Pillar 2: Mixture-Model Probability}

\textbf{Connected current modules:} Module 3, Module 8, Module 10.

\textbf{Why it matters.} This is the mathematical heart of the thesis. A proper MoE is a conditional mixture model:
\[
p(y \mid x) = \sum_k p(z=k \mid x)p(y \mid x,z=k).
\]

\textbf{Previous-knowledge bridge.} This extends ordinary probability and Bayes rule. The physics analogy is summing over latent states: the final prediction marginalizes over hidden expert assignments.

\textbf{Readings by section number.}
\begin{itemize}
  \item Bishop Ch 2.1-2.3.
  \item Bishop Ch 9.1-9.4.
  \item ESL Ch 8.5.
\end{itemize}

\textbf{Focused exercise set.}
\begin{itemize}
  \item Bishop Ch 2: Exercise 2.20.
  \item Bishop Ch 9: Exercises 9.7, 9.12, and 9.15.
  \item ESL Ch 8: Exercise 8.1 or 8.2.
  \item Applied: implement a 2-component GMM and compare with sklearn.
  \item Applied: implement a 3-expert soft MoE on synthetic regime data.
\end{itemize}

\textbf{Thesis deliverable.} You should be able to derive
\[
\mathcal{L}(\theta) =
-\sum_i \log \sum_k \pi_k(x_i;\theta)p(y_i \mid x_i,z_i=k;\theta)
\]
and explain EM, gradient descent, hard routing, and soft routing.

\section{Pillar 3: Information Theory for Gates}

\textbf{Connected current modules:} Module 3, Module 10.

\textbf{Why it matters.} You need enough information theory to talk about gate entropy, cross-entropy loss, KL divergence, expert collapse, and whether the gate uses all experts.

\textbf{Previous-knowledge bridge.} Entropy connects naturally to statistical mechanics. Softmax can be read as a Boltzmann-like distribution over experts.

\textbf{Readings by section number.}
\begin{itemize}
  \item Bishop Ch 1.6.
  \item Goodfellow Ch 3.13.
  \item Shazeer et al. (2017), Sec 4.
\end{itemize}

\textbf{Focused exercise set.}
\begin{itemize}
  \item Bishop Ch 1: Exercises 1.10 and 1.11.
  \item Applied: compute entropy of MoE gate probabilities over a batch.
  \item Applied: plot gate entropy over training and identify collapse.
  \item Applied: add an entropy penalty and compare it with Shazeer load balancing.
\end{itemize}

\textbf{Thesis deliverable.} You should be able to report
\[
H(\pi(x)) = -\sum_k \pi_k(x)\log \pi_k(x)
\]
and interpret low entropy, high entropy, and healthy expert specialization.

\section{Pillar 4: Optimization and Training Dynamics}

\textbf{Connected current modules:} Module 4, Module 6, Module 10.

\textbf{Why it matters.} MoE success often depends less on architectural elegance and more on training dynamics. Gates can collapse; experts can fail to specialize; hard routing can block gradient flow.

\textbf{Previous-knowledge bridge.} A training objective is an energy functional; regularization and load balancing are penalty terms; optimization searches for parameters that reduce the objective.

\textbf{Readings by section number.}
\begin{itemize}
  \item Prince Ch 6.1-6.5 and Ch 7.1-7.5.
  \item Goodfellow Ch 8.1-8.5.
  \item Goodfellow Ch 7.1-7.12.
\end{itemize}

\textbf{Focused exercise set.}
\begin{itemize}
  \item Prince Ch 7: Problems 7.1 and 7.2.
  \item Goodfellow Ch 6: Exercises 6.1 and 6.2.
  \item Applied: derive backprop for a 2-layer MLP and verify with \texttt{torch.autograd}.
  \item Applied: train the same MoE with and without load balancing and compare expert utilization.
\end{itemize}

\textbf{Thesis deliverable.} Explain why hard argmax routing is problematic for gradient flow, why soft routing trains end-to-end, why MoE loss is nonconvex, and why auxiliary losses are penalty methods.

\section{Pillar 5: Load Balancing and Expert Collapse}

\textbf{Connected current modules:} Module 10, Thesis Research Design.

\textbf{Why it matters.} Without balancing, the gate can route most samples to one expert. Then the model becomes a single predictor with extra unused parts.

\textbf{Previous-knowledge bridge.} This is a constrained optimization idea: low prediction loss plus a non-degenerate allocation of samples across experts. The auxiliary loss acts like a Lagrange-style penalty.

\textbf{Readings by section number.}
\begin{itemize}
  \item Shazeer et al. (2017), Sec 4.
  \item Fedus, Zoph, Shazeer (2022), Sec 2.2-2.3.
  \item Goodfellow Ch 7.1-7.2.
\end{itemize}

\textbf{Focused exercise set.}
\begin{itemize}
  \item Paper exercise: reproduce Shazeer's auxiliary load-balancing loss from Sec 4.
  \item Paper exercise: explain why expert utilization and average gate probability both matter.
  \item Applied: implement the load-balancing term.
  \item Applied: run no balancing, weak balancing, and strong balancing experiments.
\end{itemize}

\textbf{Thesis deliverable.} Include a table varying the balancing coefficient and reporting rank-IC, expert entropy, expert utilization, and turnover.

\section{Pillar 6: Dynamic Programming, HMMs, and Hamilton}

\textbf{Connected current modules:} Module 9, Module 11.

\textbf{Why it matters.} The HMM/Hamilton baseline is the classical finance comparison. You need enough dynamic programming to understand latent regime inference.

\textbf{Previous-knowledge bridge.} Think of this as a discrete-time state evolution problem. Unlike deterministic dynamical systems, the state is hidden and probabilistic. The forward algorithm propagates a probability vector through time.

\textbf{Readings by section number.}
\begin{itemize}
  \item Bishop Ch 13.1-13.3.
  \item Tsay Ch 4.1-4.6.
  \item Hamilton (1989), Sec 1-4.
\end{itemize}

\textbf{Focused exercise set.}
\begin{itemize}
  \item Bishop Ch 13: Exercises 13.1, 13.3, and 13.5.
  \item Tsay Ch 4: two end-of-chapter exercises tied to Markov switching or volatility regimes.
  \item Applied: implement the HMM forward algorithm.
  \item Applied: fit a two-state model to SPY returns and compare states with realized volatility.
\end{itemize}

\textbf{Thesis deliverable.} Explain Hamilton filtering \(p(s_t \mid y_1,\ldots,y_t)\) and contrast it with MoE gating \(p(z=k \mid x_t)\). The key distinction is temporal Markov persistence versus feature-conditioned routing.

\section{Pillar 7: Financial Validation Statistics}

\textbf{Connected current modules:} Module 5, Module 11, Module 12.

\textbf{Why it matters.} A complicated MoE is not convincing if the validation leaks information.

\textbf{Previous-knowledge bridge.} This is applied statistics: sampling error, dependence, multiple testing, and estimating uncertainty under non-i.i.d. data.

\textbf{Readings by section number.}
\begin{itemize}
  \item ESL Ch 7.1-7.12.
  \item Lopez de Prado AFML Ch 7.1-7.6.
  \item Bailey and Lopez de Prado (2014), Sec 1-4.
\end{itemize}

\textbf{Focused exercise set.}
\begin{itemize}
  \item ESL Ch 7: Exercises 7.1, 7.2, and 7.10.
  \item Applied: implement purging and embargoing for a forward-return label.
  \item Applied: compute daily rank-IC and IC information ratio.
  \item Applied: simulate many random strategies and show why the best Sharpe is biased upward.
\end{itemize}

\textbf{Thesis deliverable.} Justify chronological splits, purging and embargoing, rank-IC rather than only MSE, transaction costs, and multiple-testing caution.

\section{Optional Math: Bias-Variance for Experts and Ensembles}

\textbf{Verdict.} Learn standard ESL bias-variance before the proposal. Save non-independent ensemble bias-variance for after the proposal, and only include it if it helps the thesis discussion.

\textbf{Readings by section number.}
\begin{itemize}
  \item ESL Ch 2.9.
  \item ESL Ch 8.7 if you want ensemble context.
\end{itemize}

\textbf{Focused exercise set.}
\begin{itemize}
  \item ESL Ch 2: Exercise 2.9 or one nearby bias-variance exercise from your edition.
  \item Applied: train several experts with different seeds and measure prediction correlation.
  \item Applied: compare an average ensemble to a gated MoE.
\end{itemize}

\textbf{Thesis deliverable.} Use this to support a short discussion: MoE is not just an ensemble; it is a conditional ensemble whose weights depend on \(x\).

\section{Before-Proposal Math Plan}

\subsection{Block A: NEC Probability and Gate Math}

ESL Ch 2.1-2.9; Bishop Ch 2.1-2.3; Bishop Ch 9.1-9.4; ESL Ch 8.5. Implement GMM and soft MoE on synthetic data.

\subsection{Block B: Optimization and Deep Learning}

Prince Ch 6.1-6.5 and 7.1-7.5; Goodfellow Ch 8.1-8.5. Derive 2-layer backprop, implement a PyTorch MLP, and verify gradients.

\subsection{Block C: HMM / Hamilton Baseline}

Bishop Ch 13.1-13.3; Tsay Ch 4.6; Hamilton (1989), Sec 1-4. Implement the HMM forward algorithm and fit a two-state volatility model to SPY.

\subsection{Block D: Financial Validation}

ESL Ch 7.1-7.12; AFML Ch 7.1-7.6. Implement purged walk-forward splitting, rank-IC, and a decile backtest.

\subsection{Proposal-Level Math Deliverables}

Before proposal, be able to write and defend:

\begin{enumerate}
  \item MoE conditional mixture equation.
  \item Original NEC hard-routing defect.
  \item Corrected soft-routing equation.
  \item Load-balancing loss at a high level.
  \item HMM/Hamilton baseline at a high level.
  \item Walk-forward validation and purging logic.
\end{enumerate}

\section{After-Proposal Math Plan}

This is where the deeper theory belongs.

\textbf{Deepen mixture-model theory:} responsibilities, EM versus gradient descent, Gaussian mixture likelihood, mixture density network connection, identifiability, and label switching.

\textbf{Deepen optimization and load balancing:} expert collapse dynamics, gate entropy, Shazeer versus Switch load balancing, penalty-method interpretation, and balancing coefficient sensitivity.

\textbf{Deepen HMM and regime-switching:} forward-backward smoothing, Viterbi path, Baum-Welch updates, Markov-switching regression, regime persistence, and transition matrix interpretation.

\textbf{Deepen financial statistics:} deflated Sharpe, multiple testing, confidence intervals for IC, robustness across market regimes, and transaction-cost sensitivity.

\subsection{Thesis-Level Math Deliverables}

By January 2027, the thesis should include:

\begin{enumerate}
  \item A formal MoE model section.
  \item A short EM/HMM comparison section.
  \item A corrected NEC architecture section.
  \item A load-balancing and gate-diagnostics section.
  \item A validation/statistical testing section.
\end{enumerate}

\section{What Not To Overlearn}

Do not spend major time on these unless your supervisor specifically pushes you there:

\begin{itemize}
  \item stochastic calculus,
  \item measure-theoretic probability,
  \item PDE theory,
  \item advanced random matrix theory,
  \item optimal transport,
  \item PAC-Bayes,
  \item deep information theory.
\end{itemize}

\section{Final Math Thesis Spine}

The thesis math can be organized as:

\begin{enumerate}
  \item Conditional mixture probability: \(p(y \mid x)=\sum_k p(z=k \mid x)p(y \mid x,z=k)\).
  \item NEC defect: three-class gate -> binary hard route.
  \item Corrected MoE: \(\pi_k(x)=\mathrm{softmax}(g_k(x))\), \(\hat{y}=\sum_k \pi_k(x)f_k(x)\).
  \item Training objective: prediction loss plus \(\alpha\) times load-balancing loss.
  \item Classical baseline: HMM / Hamilton filter with Markov state persistence.
  \item Financial evaluation: walk-forward validation, purging, rank-IC, transaction-cost backtest.
\end{enumerate}

If you can explain and implement those six pieces, the math foundation is strong enough for the proposal. The deeper derivations can be added during the thesis phase.

\end{document}
